% Gemini theme
% https://github.com/anishathalye/gemini
%
% We try to keep this Overleaf template in sync with the canonical source on
% GitHub, but it's recommended that you obtain the template directly from
% GitHub to ensure that you are using the latest version.

\documentclass[final]{beamer}

% ====================
% Table
% ====================

\usepackage[table,dvipsnames]{xcolor}
\usepackage{siunitx}

\definecolor{scorebad}{HTML}{D55E00}
\definecolor{scoremidbad}{HTML}{E69F00}
\definecolor{scoremidgood}{HTML}{F0E442}
\definecolor{scoregood}{HTML}{009E73}
\definecolor{softgray}{RGB}{245,247,248}
\definecolor{linegray}{RGB}{210,216,220}

\sisetup{round-mode=places,round-precision=3,detect-weight=true,detect-family=true}
\newcommand{\wrapcell}[2]{\if\relax\detokenize{#1}\relax #2\else #1{#2}\fi}
\newcommand{\scorecell}[2][]{%
  \begingroup\setlength{\fboxsep}{1.5pt}%
  \ifdim #2 pt<0.5pt
    \cellcolor{scorebad!35}{\wrapcell{#1}{\num{#2}}}%
  \else\ifdim #2 pt<0.75pt
    \cellcolor{scoremidbad!35}{\wrapcell{#1}{\num{#2}}}%
  \else\ifdim #2 pt<0.85pt
    \cellcolor{scoremidgood!35}{\wrapcell{#1}{\num{#2}}}%
  \else
    \cellcolor{scoregood!35}{\wrapcell{#1}{\num{#2}}}%
  \fi\fi\fi
    \endgroup
}

% ====================
% Math
% ====================

\DeclareMathOperator*{\argmax}{arg\,max}
\DeclareMathOperator*{\argmin}{arg\,min}
\DeclareMathOperator{\proj}{proj}
\newcommand{\KL}[2]{\operatorname{KL}\!\left(#1\,\Vert\,#2\right)}
\newcommand{\X}{\mathcal{X}}
\newcommand{\Pd}{\mathcal{P}}
\newcommand{\Sd}{\mathbb{S}}
\newcommand{\markov}{\mathcal{M}}
\newcommand{\reciprocal}{\mathcal{R}^{\mathrm{ref}}}
\newcommand{\process}{\mathcal{P}}


% ====================
% Logo Tags
% ====================

\newcommand{\instlogo}[3][3.5cm]{%
  \begin{tikzpicture}[baseline=(img.base)]
    \node (img) {\includegraphics[height=#1]{#2}};
    \node[
      overlay,
      anchor=north east,
      xshift=8mm,
      yshift=2mm,
      text=black,
      font=\large,
    ] at (img.north east) {#3};
  \end{tikzpicture}%
}

% ====================
% Packages
% ====================

\usepackage[T1]{fontenc}
\usepackage{lmodern}
\usepackage[size=custom,width=164,height=82,scale=1.15]{beamerposter}
\usetheme{gemini}
\usecolortheme{mit}
\usepackage{graphicx}
\usepackage{booktabs}
\usepackage{tikz}
\usepackage{pgfplots}
\pgfplotsset{compat=1.14}
\usepackage{anyfontsize}
\usepackage[most]{tcolorbox}
\usetikzlibrary{arrows.meta}
\usepackage{subcaption}
\usepackage{multirow}
\usepackage{adjustbox}

% ====================
% Lengths
% ====================

% If you have N columns, choose \sepwidth and \colwidth such that
% (N+1)*\sepwidth + N*\colwidth = \paperwidth
\newlength{\sepwidth}
\newlength{\colwidth}
\setlength{\sepwidth}{0.025\paperwidth}
\setlength{\colwidth}{0.3\paperwidth}

\newcommand{\separatorcolumn}{\begin{column}{\sepwidth}\end{column}}

% ====================
% Title
% ====================

\title{Entering the Era of Discrete Diffusion Models: \\ A Benchmark for Schr\"odinger Bridges and Entropic Optimal Transport}

\author{Xavier Aramayo Carrasco \inst{1} \and Grigoriy Ksenofontov\inst{1}\inst{2} \and Aleksei Leonov\inst{2}\inst{3} \and Iaroslav Koshelev\inst{3} \and Alexander Korotin\inst{1}\inst{4}}

\institute[shortinst]{\inst{1}Applied AI Institute \samelineand \inst{2}Moscow Independent Research Institute of Artificial Intelligence \samelineand \inst{3}AI Foundation and Algorithm Lab \samelineand \inst{4}AXXX}

% ====================
% Footer (optional)
% ====================

\footercontent{
  \href{https://iclr.cc/}{https://iclr.cc/} \hfill
  ICLR 2026, Rio de Janeiro, Brazil \hfill
  \href{mailto:gregkseno@gmail.com}{gregkseno@gmail.com}
}
% (can be left out to remove footer)

% ====================
% Logo (optional)
% ====================

% use this to include logos on the left and/or right side of the header:
\logoleft{
    \includegraphics[height=8cm]{images/github_qr.png}
    \includegraphics[height=8cm]{images/hf_qr.png}
    \includegraphics[height=8cm]{images/arxiv_qr.png}
    
}
\logoright{
  \shortstack[c]{%
    \instlogo{images/AppliedAI_logo.png}{1}\hspace{10mm}%
    \instlogo{images/axxx_logo.png}{4}%
    \\[10mm]
    \instlogo{images/mirai_logo.png}{2}%
  }
}

% ====================
% Body
% ====================

\begin{document}

\begin{frame}[t]
\begin{columns}[t]
\separatorcolumn

\begin{column}{\colwidth}

  \begin{block}{Background: Entropic Optimal Transport and Schr\"odinger Bridge Problems}
      
        \vspace{0.5cm}
        \underline{\textbf{Problem setting.}} Given marginals $p_0$, $p_1$ on a discrete space $\X$ and a reference process $q^{\mathrm{ref}}$, the dynamic and static formulations are equivalent views of the same entropic transport problem.
        
        \vspace{3mm}
        \begin{minipage}[t]{0.48\linewidth}
            \begin{center}
                \textbf{Dynamic Schr\"odinger Bridge (SB)}
            \end{center}
            Seeks a stochastic process $\color{blue}{q}(x_0,x_{\mathrm{in}},x_1)$:
            \vspace{1.5cm}
            \[
            \min_{{\color{blue}{q}}\in \Pi_N(p_0,p_1)} \KL{{\color{blue}{q}}(x_0,x_{\mathrm{in}},x_1)}{q^{\mathrm{ref}}(x_0,x_{\mathrm{in}},x_1)}.
            \]
            
            \vspace{1.2cm}
            $\Pi_N(p_0,p_1)$: processes with $p_0$ at $t\!=\!0$ and $p_1$ at $t\!=\!1$.
            
            \vspace{2mm}
            \begin{figure}
                \centering
                \hspace{-10cm}
                \includegraphics[width=0.5\linewidth]{images/csb.png}
            \end{figure}
        \end{minipage}
        \hfill
        \vrule
        \hfill
        \begin{minipage}[t]{0.48\linewidth}
            \begin{center}
                \textbf{Static SB / Entropic Optimal Transport (EOT)}
            \end{center}
            Seeks a joint distribution $\color{blue}{q}(x_0,x_1)$:
            \[
            \min_{{\color{blue}q}\in \Pi(p_0,p_1)} \operatorname{KL}\bigl({\color{blue}q}(x_0,x_1)\Vert q^{\mathrm{ref}}(x_0,x_1)\bigr),
            \]
            \[
            \min_{{\color{blue}q}\in \Pi(p_0,p_1)} \mathbb{E}_{{\color{blue}q}}[-\log q^{\mathrm{ref}}(x_1|x_0)]-H({\color{blue}q}).
            \]
            $\Pi(p_0,p_1)$: joint distributions with marginals $p_0,p_1$.
            
            \vspace{2mm}
            \begin{figure}
                \centering
                \hspace{-10cm}
                \includegraphics[width=0.5\linewidth]{images/ceot.png}
            \end{figure}
        \end{minipage}
  \end{block}

    \begin{alertblock}{Motivation: Discrete and Continuous spaces}
        EOT/SB is well studied in continuous spaces ($\X=\mathbb{R}^D$), but remains underexplored in discrete spaces ($\X=\Sd^D$, $\Sd=\{0,\ldots,S-1\}$): only CSBM is currently available, and no standard evaluation pipeline exists.

        \vspace{0.5cm}
        \begin{center}
            \setlength{\fboxsep}{12pt}
            \fbox{\parbox{0.9\linewidth}{
                \centering \underline{\textbf{Goal.}} \textbf{Establish the first benchmark} for discrete-space EOT/SB solvers and \textbf{expand the family of available solvers}.
            }}
        \end{center}
    
    \end{alertblock}


    \begin{block}{Benchmark: Reverse-Engineering Ground Truth}
        \underline{\textbf{Key idea.}} Instead of fixing $(p_0,p_1)$ and solving EOT/SB directly, we choose $p_0$ and construct ${\color{blue}q^*}(x_1 \mid x_0)$. This defines the joint
        ${\color{blue}q^*}(x_0,x_1)=p_0(x_0)\,{\color{blue}q^*}(x_1 \mid x_0)$,
        and hence induces $p_1$ through the following theorem.

        \vspace{1cm}
        \begin{tcolorbox}[colback=softgray,colframe=linegray,boxrule=1pt,arc=4mm]
            \textbf{Theorem (Benchmark Pair Construction for Discrete-Space EOT/SB).}
            Let $p_0\in\Pd(\X)$ and $v^*:\X\to\mathbb{R}$. Consider a joint distribution ${\color{blue}q^*}\in\Pd(\X^2)$ such that ${\color{blue}q^*}(x_0)=p_0(x_0)$ and 
            $$
                {\color{blue}q^*}(x_1|x_0) \propto v^*(x_1)\,q^{\mathrm{ref}}(x_1|x_0).
            $$
            If $p_1(x_1):={\color{blue}q^*}(x_1)$ is its second marginal, then ${\color{blue}q^*}$ with $q^{\mathrm{ref}}$ defines the discrete-space EOT/SB between $p_0$ and $p_1$.
        \end{tcolorbox}
    \end{block}
        
    \begin{block}{Benchmark: Key Ingredients of Tractability}
        \begin{minipage}[t]{0.48\linewidth}
            We parametrize $v^*$ via \textbf{the Canonical Polyadic (CP) decomposition}:
            \begin{equation*}
                v^{*}(x_1)=\sum_{k=1}^{K}\beta_k\prod_{d=1}^{D} r_k^d[x_1^d].
            \end{equation*}
    
            \textbf{\color{Green}Universal approximation:} On a finite space, CP approximates $v^*(x_1)$ arbitrarily well as $K$ grows.
    
        \end{minipage}
        \hfill
        \vrule
        \hfill
        \begin{minipage}[t]{0.48\linewidth}
            We use \textbf{factorizable} reference processes $q^{\text{ref}}$:
            \begin{equation*}
                q^{\textup{ref}}(x_1 | x_0)=\prod_{d=1}^D q^{\textup{ref}}(x_1^d | x_0^d).
            \end{equation*} 
    
            \vspace{-4mm}
            
            \begin{itemize}\setlength{\itemindent}{0pt}
                \item \textbf{Gaussian $q^{\textup{ref}}$:} favors nearby transitions; used for ordered spaces.
                \item \textbf{Uniform $q^{\textup{ref}}$:} assigns equal mass; used when no structure is assumed.
            \end{itemize}

        \end{minipage}
    \end{block}
\end{column}

\separatorcolumn

\begin{column}{\colwidth}

    \begin{exampleblock}{Benchmark: Resulting Parameterization}
        \begin{minipage}[t]{0.48\linewidth}
            \begin{center}
                \textbf{Conditional (one-step) distribution}
            \end{center}
            For static SB and EOT problems, we have:
            \begin{gather*}
                {\color{blue}q^*}(x_1 | x_0) = \frac{1}{c^*(x_0)} \sum_{k=1}^{K} \beta_k \prod_{d=1}^{D} \Big[r_k^d[x_1^d] \, q^{\textup{ref}}(x_1^d | x_0^d)\Big],
                \\
                c^*(x_0) = \sum_{k=1}^{K} \beta_k \prod_{d=1}^{D} \Big(\sum_{x_1^d} r_k^d[x_1^d] \, q^{\textup{ref}}(x_1^d | x_0^d)\Big).
            \end{gather*}    
        \end{minipage}
        \hfill
        \vrule
        \hfill
        \begin{minipage}[t]{0.48\linewidth}
            \begin{center}
                \textbf{Transition (multi-step) distributions}
            \end{center}
            For the dynamic SB problem, we have:
            \vspace{4mm}
            \begin{gather*}
                {\color{blue}q^*}(x_{t_n} | x_{t_{n-1}}) \propto 
                q^{\textup{ref}}(x_{t_n} | x_{t_{n-1}})
                \sum_{k=1}^{K} \beta_k \prod_{d=1}^{D} u_{k,t_n}^d[x_{t_n}^d],
                \\
                u^d_{k,t_n}[x^d] = \sum_{x_1^d} q^{\textup{ref}}(x_1^d | x^d) \, r_k^d[x_1^d].
            \end{gather*}
        \end{minipage}
        \vspace{-2mm}
        \begin{center}
            \begin{tikzpicture}[>=stealth]
            
                % Images
                \node (img0) at (0,0)
                    {\includegraphics[width=0.17\linewidth]{images/p0.png}};
                \node (img1) at (11.4,0)
                    {\includegraphics[width=0.17\linewidth]{images/t1.png}};
                \node (img2) at (22.8,0)
                    {\includegraphics[width=0.17\linewidth]{images/t3.png}};
                \node (img3) at (34.2,0)
                    {\includegraphics[width=0.176\linewidth]{images/t_final.png}};
                
                % Small arrows between images
                \draw[-{Stealth[length=10pt,width=10pt]}, very thick, shorten <=8pt, shorten >=8pt]
                    (img0.east) -- (img1.west)
                    node[midway, above=6pt, align=center]
                    {\scriptsize ${\color{blue}q^*}(x_{t_1}\mid x_0)$};
                
                \draw[-{Stealth[length=10pt,width=10pt]}, very thick, shorten <=8pt, shorten >=8pt]
                    (img1.east) -- (img2.west)
                    node[midway, above=6pt, align=center]
                    {\scriptsize ${\color{blue}q^*}(x_{t_2}\mid x_{t_1})$};
                
                \draw[-{Stealth[length=10pt,width=10pt]}, very thick, shorten <=8pt, shorten >=8pt]
                    (img2.east) -- (img3.west)
                    node[midway, above=6pt, align=center]
                    {\scriptsize ${\color{blue}q^*}(x_1\mid x_{t_2})$};
                
                % Bottom arrow
                \draw[-{Stealth[length=10pt,width=10pt]}, very thick]
                    ([yshift=-6pt]img0.south)
                    -- ([yshift=-28pt]img0.south)
                    -- node[midway, below=6pt]
                       {\scriptsize ${\color{blue}q^*}(x_1\mid x_0)$}
                       ([yshift=-28pt]img3.south)
                    -- ([yshift=-6pt]img3.south);
            
            \end{tikzpicture}
        \end{center}
    \end{exampleblock}
    
   % \begin{block}{High-Dimensional Gaussian Mixture Benchmark Construction}
   %      \begin{figure}
   %      \centering
   %          \begin{minipage}[t]{0.42\linewidth}
   %              \centering
   %              \vspace{-9.5cm}
   %              \begin{tabular}{rc}
   %                  \toprule
   %                  \textbf{Parameter} & \textbf{Value} \\
   %                  \midrule
   %                  Number of dimensions ($D$) & $\{2,16,64\}$ \\
   %                  Number of categories ($S$) & $50$ \\
   %                  Number of CP components ($K$) & $5$ \\
   %                  Gaussian reference ($\gamma$) & $\{0.02,0.05\}$ \\
   %                  Uniform reference ($\gamma$) & $\{0.005,0.01\}$ \\
   %                  Number of timesteps ($N+1$) & $128$ \\
   %                  \bottomrule
   %              \end{tabular}
   %          \end{minipage}
   %          \begin{subfigure}[t]{0.25\linewidth}
   %              \centering
   %              \includegraphics[width=0.8\linewidth]{images/benchmark/d2_g002.png}
   %              \caption{\centering Gaussian with $\gamma\!=\!0.02$}
   %          \end{subfigure}
   %          \begin{subfigure}[t]{0.25\linewidth}
   %              \centering
   %              \includegraphics[width=0.8\linewidth]{images/benchmark/d2_u0005.png}
   %              \caption{\centering Uniform with $\gamma\!=\!0.005$}
   %         \end{subfigure}
   %      \end{figure}
   %  \end{block}

    \begin{block}{Discrete-State EOT/SB solvers}
        \vspace{-0.3cm}
        \begin{minipage}[t]{0.48\linewidth}
            \begin{center}
               \underline{\textbf{\textit{PRIOR}: Categorical Schr\"odinger}} \\ \underline{\textbf{Bridge Matching (CSBM)}}
            \end{center}

            \textbf{Idea.} Alternating projections onto the Markov $\mathcal{M}$ and the reciprocal $\mathcal{R}^{\mathrm{ref}}$ sets:
            \[
                q^{2l+1}=\proj_{\mathcal M}(q^{2l}),\qquad
                q^{2l+2}=\proj_{\mathcal R^{\mathrm{ref}}}(q^{2l+1}).
            \]
                
            \textbf{Training.} Bridge matching (BM) optimization:
            \vspace{-5mm}
            \begin{multline*}
                \argmin_{\color{blue}\theta}\mathbb{E}_{q^{2l+1}(x_0,x_1)}
                \Bigg[
                \sum_{n=1}^{N}
                \mathbb{E}_{q^{\mathrm{ref}}_{n-1}}
                \KL{q^{\mathrm{ref}}_{n}}{{\color{blue}q_{\theta,n}}}
                -
                \\[-10mm]
                -
                \mathbb{E}_{q^{\mathrm{ref}}_{N}}
                \log {\color{blue}q_{\theta}}(x_1\mid x_{t_N})
                \Bigg].
            \end{multline*}
            
            \centering
            \includegraphics[width=0.5\linewidth]{images/imf.png}
            \vspace{4mm}
        \end{minipage}
        \hfill
        \vrule
        \hfill
        \begin{minipage}[t]{0.48\linewidth}
            \begin{center}
                \underline{\textbf{\textit{NEW}: $\alpha$-Categorical Schr\"odinger}} \\ \underline{\textbf{Bridge Matching ($\alpha$-CSBM)}}
            \end{center}

            \textbf{Idea.} Relaxed projections with step size $\alpha$:
            \vspace{1.2cm}
            \[
                q^{l+1}=(1-\alpha)q^l+\alpha\,\proj_{\mathcal R^{\mathrm{ref}}}\!\big(\proj_{\mathcal M}(q^l)\big).
            \]

            \textbf{Training.} Joint forward/backward BM optimization:
            \vspace{8mm}
            \begin{gather*}
                \argmin_{\color{blue}\theta}\tfrac{1}{2}\Big(
                \mathrm{KL}(\overrightarrow{r_{\mathrm{sg}}}\,\|\,\overleftarrow{\color{blue}q_{\theta}})+
                \mathrm{KL}(\overleftarrow{r_{\mathrm{sg}}}\,\|\,\overrightarrow{{\color{blue}q_{\theta}}})
                \Big), 
                \\
                \text{where } r_{\mathrm{sg}}=\proj_{\mathcal{R}^{\mathrm{ref}}}({\color{blue}q_{\theta}}).
            \end{gather*}

            \vspace{5mm}
            \centering
            \includegraphics[width=0.5\linewidth]{images/alpha-imf.png}
            \vspace{4mm}
        \end{minipage}
        
        \hrule

        \vspace{-4mm}
        \begin{center}
            \underline{\textbf{\textit{NEW}: Discrete Light Schr\"odinger Bridge (Matching) (DLightSB(-M))}}
        \end{center}
        \vspace{-6mm}
        \textbf{Idea.} DLightSB(-M) parameterizes ${\color{blue}v_\theta}$ so that ${\color{blue}q_\theta}$ always remains an SB. Hence, solving SB reduces to matching the target marginal $p_1$ via {\color{Green}a single projection}.

        \vspace{-3mm}
        \begin{minipage}[t]{0.48\linewidth}
            \begin{center}
                \underline{\textbf{DLightSB}}
            \end{center}
            
            Focuses static target ${\color{blue}q_{\theta}}(x_1 | x_0)$.

            \textbf{Training.} Minimize $\KL{{\color{blue}q^*}}{{\color{blue}q_\theta}}$ via:
            \vspace{2cm}
            \[
                \argmin_{\color{blue}\theta}\mathbb{E}_{p_0(x_0)}\big[\log {\color{blue}c_{\theta}}(x_0)\big] - \mathbb{E}_{p_1(x_1)}\big[\log {\color{blue}v_{\theta}}(x_1)\big].
            \]
        \end{minipage}
        \hfill
        \vrule
        \hfill
        \begin{minipage}[t]{0.48\linewidth}
             \begin{center}
                \underline{\textbf{DLightSB-M}}
            \end{center}
            Focuses dynamic target ${\color{blue}q_{\theta}}(x_{t_n} | x_{t_{n-1}})$.
        
            \textbf{Training.} Minimize $\KL{r}{{\color{blue}q_\theta}}$ via BM optimization:
            \vspace{-1mm}
            \begin{multline*}
                \argmin_{\color{blue}\theta}\mathbb{E}_{q^{2l+1}(x_0,x_1)}
                \Big[
                \sum_{n=1}^{N}
                \mathbb{E}_{q^{\mathrm{ref}}_{n-1}}
                \KL{q^{\mathrm{ref}}_{n}}{{\color{blue}q_{\theta,n}}}
                -
                \\[-10mm]
                -
                \mathbb{E}_{q^{\mathrm{ref}}_{N}}
                \log {\color{blue}q_{\theta}}(x_1\mid x_{t_N})
                \Big].
            \end{multline*}
        \end{minipage}        
    \end{block}

\end{column}

\separatorcolumn

\begin{column}{\colwidth}

  \begin{block}{Evaluation Protocol}
        % \small
        \begin{minipage}[t]{0.39\linewidth}
            \underline{\textbf{Baselines:}}
            \begin{itemize}\setlength{\itemindent}{0pt}
              \item \textit{Independent}: sample $x_1$ directly from $p_1$; ignores the joint distribution.
              \item \textit{Reference}: sample $x_1$ from $q^{\mathrm{ref}}$; follows prior transport.
              \item \textit{Feature-wise SB}: solve 1D SB independently for each dimension; ignores cross-dimensional dependence.
            \end{itemize}
        \end{minipage}
        \hfill
        \vrule
        \hfill
        \begin{minipage}[t]{0.58\linewidth}
            \underline{\textbf{Metrics}:}
            Shape Score (SSM) and Trend Score (TSM) average over 1D and 2D marginals, respectively:
            \vspace{-5mm}
            \begin{gather*}
                \mathrm{SSM}^d = 1-\frac12\sum_{x^d=0}^{S-1}\left|\tilde q^*(x^d)-\tilde q_\theta(x^d)\right|,
                \\[-2mm]
                \mathrm{TSM}^{d_m,d_n}=1-\frac12\sum_{x^{d_m}=0}^{S-1}\sum_{x^{d_n}=0}^{S-1}\left|\tilde q^*(x^{d_m},x^{d_n})-\tilde q_\theta(x^{d_m},x^{d_n})\right|.
            \end{gather*}
            Here $\tilde q(x)=\frac1{|I|}\sum_{i\in I}\delta_{x^{(i)}}$ is the empirical distribution. Conditional variants average SSM and TSM over $x_0\sim p_0$ and multiple $x_1$ per $x_0$.
        \end{minipage}
    \end{block}

    \begin{block}{Conditional Trend Score Results}
        \vspace{-0.4cm}
        \begin{center}
            \begin{adjustbox}{width=0.99\linewidth}
            \begin{tabular}{rcccccccccccccc}
            & & & \multicolumn{4}{c}{$D=2$} & \multicolumn{4}{c}{$D=16$} & \multicolumn{4}{c}{$D=64$}\\
            \cmidrule(lr){4-7} \cmidrule(lr){8-11} \cmidrule(l){12-15}
            & & & \multicolumn{2}{c}{gaussian} & \multicolumn{2}{c}{uniform} & \multicolumn{2}{c}{gaussian} & \multicolumn{2}{c}{uniform} & \multicolumn{2}{c}{gaussian} & \multicolumn{2}{c}{uniform}\\
            \cmidrule(lr){4-5} \cmidrule(lr){6-7} \cmidrule(lr){8-9} \cmidrule(lr){10-11} \cmidrule(lr){12-13} \cmidrule(l){14-15}
            Method & Loss & $N\!+\!1$ & $0.02$ & $0.05$ & $0.005$ & $0.01$ & $0.02$ & $0.05$ & $0.005$ & $0.01$ & $0.02$ & $0.05$ & $0.005$ & $0.01$\\
            \midrule
            \textit{Independent} & -- & -- & \scorecell{0.471} & \scorecell{0.779} & \scorecell{0.521} & \scorecell{0.551} & \scorecell{0.569} & \scorecell{0.676} & \scorecell{0.369} & \scorecell{0.458} & \scorecell{0.478} & \scorecell{0.509} & \scorecell{0.351} & \scorecell{0.432}\\
            \textit{Reference} & -- & -- & \scorecell{0.085} & \scorecell{0.277} & \scorecell{0.230} & \scorecell{0.273} & \scorecell{0.111} & \scorecell{0.082} & \scorecell{0.090} & \scorecell{0.095} & \scorecell{0.195} & \scorecell{0.121} & \scorecell{0.136} & \scorecell{0.110}\\
            \textit{Feature-wise SB} & -- & -- & \scorecell{0.585} & \scorecell{0.630} & \scorecell{0.704} & \scorecell{0.738} & \scorecell{0.686} & \scorecell{0.532} & \scorecell{0.351} & \scorecell{0.373} & \scorecell{0.838} & \scorecell{0.404} & \scorecell{0.532} & \scorecell{0.456}\\
            \midrule
            DLightSB & -- & -- & \textbf{\scorecell{0.908}} & \textbf{\scorecell{0.851}} & \textbf{\scorecell{0.873}} & \textbf{\scorecell{0.872}} & \scorecell{0.763} & \textbf{\scorecell{0.838}} & \textbf{\scorecell{0.839}} & \textbf{\scorecell{0.838}} & \textbf{\scorecell{0.854}} & \textbf{\scorecell{0.747}} & \textbf{\scorecell{0.817}} & \textbf{\scorecell{0.773}}\\
            \midrule
            \multirow{4}{*}{CSBM} & \multirow{2}{*}{KL} & 16 & \scorecell{0.643} & \scorecell{0.610} & \scorecell{0.801} & \scorecell{0.808} & \scorecell{0.674} & \scorecell{0.609} & \scorecell{0.517} & \scorecell{0.513} & \scorecell{0.756} & \scorecell{0.694} & \scorecell{0.517} & \scorecell{0.537}\\
             &  & 64 & \underline{\scorecell{0.837}} & \scorecell{0.794} & \scorecell{0.840} & \scorecell{0.846} & \scorecell{0.801} & \scorecell{0.738} & \scorecell{0.572} & \scorecell{0.624} & \underline{\scorecell{0.788}} & \scorecell{0.698} & \scorecell{0.486} & \scorecell{0.536}\\
             & \multirow{2}{*}{MSE} & 16 & \scorecell{0.444} & \scorecell{0.627} & \scorecell{0.692} & \scorecell{0.728} & \scorecell{0.658} & \scorecell{0.574} & \scorecell{0.392} & \scorecell{0.418} & \scorecell{0.586} & \scorecell{0.655} & \scorecell{0.446} & \scorecell{0.463}\\
             &  & 64 & \scorecell{0.227} & \scorecell{0.748} & \scorecell{0.686} & \scorecell{0.673} & \scorecell{0.726} & \scorecell{0.716} & \scorecell{0.482} & \scorecell{0.551} & \scorecell{0.496} & \underline{\scorecell{0.714}} & \scorecell{0.479} & \scorecell{0.478}\\
            \midrule
            \multirow{4}{*}{$\alpha$-CSBM} & \multirow{2}{*}{KL} & 16 & \scorecell{0.588} & \scorecell{0.616} & \scorecell{0.819} & \scorecell{0.814} & \scorecell{0.610} & \scorecell{0.659} & \scorecell{0.545} & \scorecell{0.557} & \scorecell{0.773} & \scorecell{0.659} & \scorecell{0.484} & \scorecell{0.515}\\
             &  & 64 & \scorecell{0.734} & \scorecell{0.794} & \scorecell{0.834} & \scorecell{0.844} & \textbf{\scorecell{0.816}} & \scorecell{0.798} & \scorecell{0.563} & \scorecell{0.642} & \underline{\scorecell{0.793}} & \scorecell{0.677} & \scorecell{0.486} & \scorecell{0.514}\\
             & \multirow{2}{*}{MSE} & 16 & \scorecell{0.470} & \scorecell{0.607} & \scorecell{0.767} & \scorecell{0.781} & \scorecell{0.604} & \scorecell{0.649} & \scorecell{0.449} & \scorecell{0.512} & \scorecell{0.592} & \scorecell{0.618} & \scorecell{0.448} & \scorecell{0.457}\\
             &  & 64 & \underline{\scorecell{0.838}} & \scorecell{0.782} & \scorecell{0.758} & \scorecell{0.785} & \scorecell{0.587} & \scorecell{0.761} & \scorecell{0.486} & \scorecell{0.561} & \scorecell{0.488} & \scorecell{0.687} & \scorecell{0.466} & \scorecell{0.486}\\
            \midrule
            \multirow{4}{*}{DLightSB-M} & \multirow{2}{*}{KL} & 16 & \scorecell{0.817} & \textbf{\scorecell{0.848}} & \underline{\scorecell{0.860}} & \underline{\scorecell{0.864}} & \scorecell{0.702} & \scorecell{0.820} & \underline{\scorecell{0.831}} & \underline{\scorecell{0.820}} & \scorecell{0.731} & \scorecell{0.622} & \scorecell{0.468} & \scorecell{0.500}\\
             &  & 64 & \scorecell{0.798} & \underline{\scorecell{0.843}} & \scorecell{0.853} & \underline{\scorecell{0.856}} & \underline{\scorecell{0.791}} & \underline{\scorecell{0.825}} & \underline{\scorecell{0.825}} & \underline{\scorecell{0.816}} & \scorecell{0.727} & \underline{\scorecell{0.710}} & \scorecell{0.493} & \underline{\scorecell{0.674}}\\
             & \multirow{2}{*}{MSE} & 16 & \scorecell{0.624} & \underline{\scorecell{0.839}} & \scorecell{0.755} & \scorecell{0.838} & \scorecell{0.451} & \scorecell{0.824} & \scorecell{0.783} & \scorecell{0.785} & \scorecell{0.532} & \scorecell{0.683} & \underline{\scorecell{0.645}} & \scorecell{0.328}\\
             &  & 64 & \scorecell{0.595} & \scorecell{0.834} & \scorecell{0.716} & \scorecell{0.822} & \scorecell{0.404} & \scorecell{0.814} & \scorecell{0.786} & \scorecell{0.763} & \scorecell{0.392} & \scorecell{0.678} & \scorecell{0.327} & \scorecell{0.600}\\
            \bottomrule
            \end{tabular}
            \end{adjustbox}
        \end{center}
    \end{block}

    \begin{exampleblock}{Key Takeaways}
        \vspace{5mm}
        \noindent
        \begin{minipage}[t]{0.31\linewidth}
            \centering
            \underline{\textbf{Baselines}}
            \vspace{2mm}
            
            \begin{itemize}\setlength{\itemindent}{0pt}
                  \item \textit{Reference} {\color{Red}degrades} because $v^*$ strongly reweights the reference.
                  \item \textit{Independent} {\color{Red}worsens} at low $\gamma$.
                  \item \textit{Feature-wise SB} {\color{Red}degrades} as dimension increases.
            \end{itemize}
            \vspace{3mm}
            \setlength{\fboxsep}{15pt}
            \fbox{\parbox{0.9\linewidth}{\centering
                  These failures highlight the key challenges our benchmarks are designed to test.
            }}
        \end{minipage}
        \hspace{0.2em}   % reduced space before rule
        \vrule
        \hspace{0.2em}   % reduced space after rule
        \begin{minipage}[t]{0.31\linewidth}
            \centering
            \underline{\textbf{DLightSB-like}}
            \vspace{2mm}
            
            \begin{itemize}\setlength{\itemindent}{0pt}
                  \item DLightSB: {\color{Green}strongest} performance across all setups.
                  \item DLightSB-M: comparable results, {\color{Orange}slight drop} due to KL.
                  \item Both {\color{Red}benefit from the inductive bias} of the benchmark.
            \end{itemize}
            \vspace{4mm}
            \setlength{\fboxsep}{15pt} 
            \fbox{\parbox{0.9\linewidth}{\centering
                DLightSB(-M) can be viewed as \textbf{oracle-like} in this benchmark.
            }}
        \end{minipage}
        \hspace{0.2em}
        \vrule
        \hspace{2mm}
        \begin{minipage}[t]{0.31\linewidth}
        \centering
        \underline{\textbf{CSBM-like}}
        \vspace{2mm}
        
        \begin{itemize}\setlength{\itemindent}{0pt}
          \item CSBM and $\alpha$-CSBM perform {\color{Orange}worse} than DLightSB(-M).
          \item $\alpha$-CSBM matches CSBM quality at {\color{Green}lower computational cost}.
          \item Higher $N$ generally improves metrics.
          \item KL consistently {\color{Green}outperforms} MSE (MSE over-smooths).
        \end{itemize}
        \end{minipage}
    \end{exampleblock}

\end{column}

\separatorcolumn
\end{columns}
\end{frame}

\end{document}
